\chapter*{Preface}
\addcontentsline{toc}{chapter}{Preface}

In this book we journey into the realm of cellular automata (CA).
CA states are displayed in various ways, allowing readers of limited
mathematical training to appreciate the beauty of CA-evolved patterns.
Each chapter has an associated software module (in the Python package
\texttt{cadyn}) that generated the chapter's examples.  Readers are
encouraged to create their own examples, perform their own experiments,
but most importantly to have fun with this software.

Chapter~1 discusses why CA are of interest and to what they have been
applied, and introduces the basic terminology and definitions.

In Chapter~2 we make an excursion into the land of 0-dimensional CA.
This gives a simple introduction to fundamental concepts used in
subsequent chapters, and introduces vocabulary shared by the fields of
continuous and discrete dynamical systems.  All examples were generated
with the module \texttt{cadyn.fds}.

In Chapter~3 we explore 1-dimensional CA and exhibit connections with
1-dimensional digital signal processing.  The examples of this chapter
were created with \texttt{cadyn.ca1d}.

Chapter~4 looks at 2-dimensional CA, whose states are easily displayed
by translating each cell's state into a pixel color.  Many visually
interesting examples are shown, and connections with central topics in
image processing are discussed.  These examples were created with
\texttt{cadyn.ca2d}, with which the reader can create a virtually
unlimited supply of colorful and geometrically rich patterns.  Artists
have found this a useful pattern-generating tool.

A particular class of ecosystem simulations, using 2-dimensional CA, is
the subject of Chapter~5.  Various experiments are performed and the
resulting population dynamics discussed, using
\texttt{cadyn.ecosystem}, which lets the user set ecosystem parameters
and observe the subsequent dynamics.

Chapter~6 introduces Stephen Wolfram's classification of CA into four
qualitative classes, relates these classes to fundamental issues in the
theory of computation, and discusses various philosophical issues that
arise in the context of CA.

An extensive bibliography and related links allow readers to continue
their studies of CA in many directions.

I would like to thank Bakersfield College administrators and the Kern
Community College District Board of Trustees for granting the
sabbatical leave that made the first edition possible.
